\documentclass[10pt]{article}

\usepackage{amsmath}
\usepackage{amssymb}
\usepackage{amsthm}
\usepackage{ mathrsfs }
\usepackage{ mathtools}
\usepackage{enumerate}
\usepackage{bbm}
\usepackage{lipsum}
\usepackage{fancyhdr}
\usepackage{tikz-cd} 
\usepackage{adjustbox} 
\usepackage{stmaryrd}
\numberwithin{equation}{section}
\usetikzlibrary{arrows}
\newcommand{\midarrow}{\tikz \draw[-triangle 90] (0,0) -- +(.1,0);}
\tikzset{commutative diagrams/.cd,
mysymbol/.style={start anchor=center,end anchor=center,draw=none}
}
\newcommand\comsymb[2][\circlearrowleft]{%
  \arrow[mysymbol]{#2}[description]{#1}}

\newtheorem{theorem}{Theorem} [section] 
\newtheorem{proposition}{Proposition}[section] 
\newtheorem{definition}{Definition}[section] 
\newtheorem{lemma}{Lemma}[section] 
\newtheorem{notation}{Notation}[section] 
\newtheorem{remark}{Remark}[section] 
\newtheorem{corollary}{Corollary} [section] 
\newtheorem{terminology}{Terminology}[section] 
\newtheorem{fact}{Fact}[section] 
\newtheorem{example}{Example}[section] 
\newtheorem{claim}{Claim}
\newenvironment{claimproof}[1]{\par\noindent\underline{Proof:}\space#1}{\hfill $\blacksquare$}

\DeclareMathOperator{\tmod}{mod}
\DeclareMathOperator{\triv}{triv}
\DeclareMathOperator{\ind}{ind}
\DeclareMathOperator{\straw}{s.t.}
\DeclareMathOperator{\ev}{ev}
\DeclareMathOperator{\pr}{pr}
\DeclareMathOperator{\Aut}{Aut}
\DeclareMathOperator{\Map}{Map}
\DeclareMathOperator{\Stab}{Stab}
\DeclareMathOperator{\Ind}{Ind}
\DeclareMathOperator{\GL}{GL}
\DeclareMathOperator{\SL}{SL}
\DeclareMathOperator{\SO}{SO}
\DeclareMathOperator{\Sp}{Sp}
\DeclareMathOperator{\esssup}{esssup}
\DeclareMathOperator{\abs}{abs}
\DeclareMathOperator{\rel}{rel}
\DeclareMathOperator{\diam}{diam}
\DeclareMathOperator{\rank}{rank}
\DeclareMathOperator{\Hom}{Hom}
\DeclareMathOperator{\Homk}{Hom_{k-alg}}
\DeclareMathOperator{\Ker}{Ker}
\DeclareMathOperator{\Image}{Im}
\DeclareMathOperator{\Dom}{Dom}
\DeclareMathOperator{\grad}{grad}
\DeclareMathOperator{\divergence}{div}
\DeclareMathOperator{\rk}{rk}
\DeclareMathOperator{\Span}{Span}
\DeclareMathOperator{\mSpec}{mSpec}
\DeclareMathOperator{\MaxSpec}{MaxSpec}
\DeclareMathOperator{\interior}{int}
\DeclareMathOperator{\supp}{supp}
\DeclareMathOperator{\id}{id}
\DeclareMathOperator{\sgn}{sgn}
\DeclareMathOperator{\Mat}{Mat}
\DeclareMathOperator{\PD}{PD}
\DeclareMathOperator{\ext}{ext}
\DeclareMathOperator{\vol}{vol}
\DeclareMathOperator{\inte}{int}
\DeclareMathOperator{\diverg}{div}
\DeclareMathOperator{\curl}{curl}
\DeclareMathOperator{\image}{im}
\DeclareMathOperator{\dR}{dR}
\DeclareMathOperator{\Ob}{Ob}
\DeclareMathOperator{\Mnfd}{Mnfd}
\DeclareMathOperator{\OpEmb}{OpEmb}
\DeclareMathOperator{\Spec}{Spec}
\DeclareMathOperator{\Spa}{Spa}
\DeclareMathOperator{\Lim}{Lim}
\DeclareMathOperator{\Gal}{Gal}
\DeclareMathOperator{\rid}{rid}
\DeclareMathOperator{\Spf}{Spf}
\DeclareMathOperator{\Rig}{Rig}
\DeclareMathOperator{\Adic}{Adic}
\DeclareMathOperator{\FSch}{FSch}
\DeclareMathOperator{\Sper}{Sper}
%\newcommand*{\name}[\num_arguments][default values]{{\color{#1}\Large #2}}
\newcommand{\defeq}{\vcentcolon=}
\newcommand{\Ouv}{\mathcal{O}}
\newcommand{\norm}[1]{\Vert #1 \Vert}
\newcommand{\opNorm}[2]{\norm{#1}_{#2\to#2}}
\newcommand{\normL}[3]{\norm{#1}_{L^{#2}(#3)}}
\newcommand{\N}[0]{\mathbb{N}}
\newcommand{\m}[0]{\mathfrak{m}}
\newcommand{\R}[0]{\mathbb{R}}
\newcommand{\Z}[0]{\mathbb{Z}}
\newcommand{\C}[0]{\mathbb{C}}
\newcommand{\Q}[0]{\mathbb{Q}}
\newcommand{\F}[0]{\mathbb{F}}
\newcommand{\G}[0]{\mathbb{G}}
\newcommand{\A}[0]{\mathbb{A}}
\newcommand{\T}[0]{\mathbb{T}}
\newcommand{\U}[0]{\mathfrak{U}}
\newcommand{\Para}[0]{\mathbb{P}}
\newcommand{\M}[0]{\mathcal{M}}
\newcommand{\maniN}[0]{\mathcal{N}}
\newcommand{\cinf}[0]{\mathcal{C}^\infty}
\newcommand{\torus}[0]{\mathbb{T}}
\newcommand{\sep}[0]{\:|\:}
\newcommand{\pd}[2]{{\frac {\partial #1} {\partial #2}}}
\newcommand{\compinc}[0]{\subset \subset}
\newcommand{\sH}[0]{\mathscr{H}}
\newcommand{\ab}[1]{\vert #1 \vert}
\newcommand{\st}[0]{\:\straw\:}
\newcommand{\g}[0]{\mathfrak{g}}
\newcommand{\p}[0]{\mathfrak{p}}
\newcommand{\fib}[1]{%
  \mathbin{\mathop{\times}\limits_{#1}}%
}
\newcommand{\tens}[1]{%
  \mathbin{\mathop{\otimes}\displaylimits_{#1}}%
}

\title{Rigid Analytic Geometry Tutorial}
\author{So Murata}
\date{SoSe 26, University of Bonn}


\begin{document}
\maketitle

%4/20

\section{Week 2}
\subsection{Motivation and Obstacles}
The motivation for the course is that, given a non-archimedian field $k$, how should we define analytic varieties and presheaves of algebraic functions over $k$.

\begin{definition}[Classical Rigid Analytic Geometry]
  Intuition, over $\overline{B(0,1)}^n\subseteq K^n$, the analytic functions should be of the form of convergent power series.
  \begin{equation*}
    f = \sum_{\alpha\in\N^n}a_\alpha x^\alpha.
  \end{equation*}
  This is a local description of the functions. We denote the set of such functions 
  \begin{equation*}
    \T_n(K)\defeq K\langle x_1,x_n\rangle,
  \end{equation*}
  which we call the Tate algebra. 
\end{definition}

\begin{remark}
  There is a one-to-one correspondence between 
  \begin{equation*}
    \Sp\T_n(K)\longleftrightarrow \overline{B_{\overline{K}}(0,1)}^n/\Gal(\overline{K}/K).
  \end{equation*}
\end{remark}

We will denote $X=\Sp\T_n(K)$. Furthermore, we can define admissible opens $U\subseteq X$, analytic functions on $U$, and the presheaf of analytic functions $\Ouv_{X}$.
Over a wide range of admissible open coverings, $\Ouv_{X}$ satisfies the sheaf requirement.

\begin{equation*}
  0\to\Ouv(X)\hookrightarrow\prod_{i}\Ouv(U_i)\rightrightarrows\prod_{i,j}\Ouv(U_i\cap U_j).
\end{equation*}

We only allow certain open coverings called admissible covers. We introduce and work on a Grothendieck topology.
\begin{remark}
  There's a non-trivial sheaf of abelian groups $\mathcal{F}$ on $X$ with $G$-topology such that $\forall x\in X,\F_x=0$.
\end{remark}

To overcome the obstacle above, we will introduce extra points, namely van der Put points.
We glue $\Sp\T_n(K)/I$ along opens to get rigid $K$-spaces.


\subsection{Raynaud's Formal Schemes}

Suppose $X$ is a $\Q_p$-scheme, we want to consider the reduction $\tmod{p}$. We find a $\Z_p$-scheme $X'$ such that 
\begin{equation*}
  X'\fib{\Z_p}\Q_p\cong X.
\end{equation*}
Note that we can't guarantee that $X'$ is unique or even it exists. Take $R$ to be the valuation ring of $K$ for example $K=\Q_p,R=\Z_p$, we want to find an object the structure over $R$, such that its generic fiber is $X$.
Say $X=\Sp\T_n(K)=\Sp K\langle x_1,\cdots,x_n\rangle$, we can take 
\begin{equation*}
  Y = \Spf R\langle x_1,\cdots,x_n\rangle,
\end{equation*}
We also set,
\begin{equation*}
  Y^{\rid}=\Sp(R\langle x_1,\cdots,x_n\rangle\hat{\tens{R}}K).
\end{equation*}
If $R$ is separated $\mathfrak{a}$-adically topologized, that is for any subset $I\subseteq R$, $U$ is open if and only if there is $n\in\N$ such that $\mathfrak{a}^n\subseteq I$, and $\bigcup_{n\geq1}\mathfrak{a}^n=0$. 
Then we have,
\begin{equation*}
  R\langle x_1,\cdots,x_n\rangle = \varprojlim_{m\geq1}R/\mathfrak{a}^n[x_1,\cdots,x_n].
\end{equation*}
We have another construction of rigid spaces, namely Berkovich's theory.
\subsection{Huber's adic spaces}

Let $k$ be algebraically closed and $K/k$ is a finitely generated field extension of transcendence degree $1$. Then there is a non-singular projective curve over $k$ with 
function field $K$. 
\par\textbf{Construction}
Take $C_K$ to be the set of discrete valuation rings in $K/k$. Take closed subsets to be the whole space and finite sets.
On each open $U$, define 
\begin{equation*}
  \Ouv(U) = \bigcap_{p\in U}R_p,
\end{equation*}
then $(C_K,\Ouv)$ is isomorphic to a non-singular curve with function field $K$.
\par Given a topological ring $A$ and and its subring $A^\triangleleft$ or sometimes denoted $A^+$(integrally closed, open, power bounded). Huber defines,
\begin{equation*}
  \Spa(A,A^\triangleleft) = \{\text{continuous valuations $v:A\to \Gamma\cup\{0\}$ such that $\forall a\in A^\triangleleft, v(a)\leq1$}\}/\sim.
\end{equation*}
\begin{remark}  There is a fully-faithful functor 
  \begin{equation*}
    r_K:\Rig(K)\to \Adic/\Spa(K,K^\circ),
  \end{equation*}
  from the category of rigid $K$-spaces, to the category of adic spaces over $\Spa(K,K^\circ)$.
\end{remark}
\begin{remark}  There is a fully-faithful functor 
  \begin{equation*}
    t:\FSch^{\text{ln}}\to \Adic/\Spa(K,K^\circ),
  \end{equation*}
  from the category of locally Noetherian formal schemes, to the category of adic spaces over $\Spa(K,K^\circ)$.
\end{remark}
\begin{remark}
  \label{remark_2_classical_modern_intersection}
  If $X$ is a qcqs rigid $K$-space, the set of van der Put points of $X$ is one-to-one correspondent to points of $r_K(X)$.
\end{remark}
The problem session goes parallel to the lecture, meaning the lecture will focus on the modern treatment while the tutorial will focus on the classical one. Remark \ref{remark_2_classical_modern_intersection} will be the meeting point of these two.
\section{Week1}
\par\textbf{Problem 2}
Let $f,g\in\T_m$. Set
\begin{equation*}
  f = \sum_{\alpha}a_\alpha x^\alpha,g=\sum_\beta b_\beta x^\beta.
\end{equation*}
Clearly,
\begin{align*}
  \norm{f+g} & = \max_{\alpha} \vert a_\alpha+b_\alpha\vert,\\
  & \leq \max_\alpha\{\vert a_\alpha\vert,\vert b_\alpha\vert\},\\
  & = \max\{\norm{f},\norm{g}\}.
\end{align*}
For $k\in K$, we have,
\begin{equation*}
  \norm{k\cdot f} = \max_\alpha\vert k a_\alpha\vert  = \vert k\vert \max_\alpha\vert a_\alpha\vert.
\end{equation*}
For the multiplicity, assume $\norm{f},\norm{g} = 1$. This means there are $\alpha_0,\beta_0$ such that 
\begin{equation*}
  a_{\alpha_0} = 1,b_{\beta_0} = 1.
\end{equation*}

Recall that $K^\circ$ is a local with maximal ideal $K^{\circ\circ}$ so set $k=K^\circ/K^{\circ\circ}$. Let $\overline{f},\overline{g}$ be reductions of some $f,g$ in $k[x_1,\cdots,x_n]$. Then $\overline{f},\overline{g}$ are non-zero so $\overline{f}\overline{g}$ is non-zero.
This shows that some coefficients of $fg$ is in $K^\circ\backslash K^{\circ\circ}$. Thus $\norm{fg}\geq1$. 
\par We have,
\begin{align*}
  \norm{fg}\leq\max_{\gamma}\max_{\alpha+\beta=\gamma}\vert a_\alpha b_\beta\vert\leq \norm{f}\norm{g}.
\end{align*}
\begin{definition}[Topological and Norm-free definition of boundedness, Power boundedness]
  Let $A$ be a topological ring and a subset $X$ of $A$ is called bounded if for any neighbourhood $U$ of $0$, there is a open neighbourhood $V$ of $0$ such that 
  \begin{equation*}
    V\cdot X \subseteq U.
  \end{equation*}
    An element $a\in A$ is called power-bounded if $\{a^n\sep n\geq1\}$ is bounded.
\end{definition}
\begin{example}
  Consider $A=k[x]/(x^2)$. 
\end{example}
\begin{remark}
  If we consider the quotient topology on $A/I$ for $I$ open, then $A/I$ will be equipped with discrete topology. This is because $a+I$ is open thus the reduction $\{\overline{a}\}$ is open.
\end{remark}
If $A$ is a $K$-Banach algebra and the topological structure of $K$ is not discrete by definition. An open ideal $I$ in $A$ will make $K$-Banach algebra $A/I$ discrete, in particular, $K\subseteq A/I$ is discrete thus $A/I$ is trivial.
\par Let $R=K^\circ$ be the valuation ring of $K$. We then define,
\begin{equation*}
  R\langle x_1,\cdots,x_n\rangle = \left\{f=\sum_{\alpha\in\N^m}a_\alpha x^\alpha\sep \vert a_\alpha\vert \to 0\right\}.
\end{equation*}
Define the topology on $R\langle x_1,\cdots,x_n\rangle$ as the subset topology of $\T_m$. Then this admits a non-trivial open proper ideal. We take a pseudo-uniformizer $\varpi\in R^{\circ\circ}=K^{\circ\circ}$ (i.e. $\vert\varpi\vert<1$), then $(\varpi)$ is an open ideal of $A$. Moreover,
\begin{equation*}
  \{(\varpi^n)\}_{n\in\N}=\left\{\sum_{\alpha\in\N^m}a_\alpha x^\alpha\sep \vert a_\alpha\vert\to 0 ,\vert a_\alpha\vert \leq\vert \varpi^n\vert \right\}.
\end{equation*} 
is a basis of neighborhoods of $0$. 
\par What is the completion of $R[x]$ with respect to $(\varpi)$?
\begin{definition}
  The completion of $R[x]$ with respect to $(\varpi)$ is defined as 
  \begin{equation*}
    \varprojlim_n R[x]/(\varpi)^n = R\langle x\rangle, \varprojlim_n R[x]/x^n = R\llbracket x\rrbracket,\varprojlim_n R[x]/(x,\varpi)^n = R\llbracket x\rrbracket.
  \end{equation*}
\end{definition}

Question is 
\begin{equation*}
  R\langle x_1,\cdots, x_n\rangle\tens{R}k \cong \T_m
\end{equation*}
and 
\begin{equation*}
  K\langle x\rangle\tens{K}K\langle y\rangle \cong K\langle x,y\rangle
\end{equation*}
hold? It turns out the first one does with $K=R\left[{\frac 1 \varpi}\right]$. However, the second one doesnt hold. But we have 
\begin{equation*}
  K\langle x\rangle\tens{K} K\langle y\rangle
\end{equation*}
 is dense in $K\langle x,y\rangle$.
 \par Suppose $K=\overline{K}$ from now on. Here $\overline{B(0,1)}\subset K^n$ is endowed with maximum norm. We denote,
 \begin{equation*}
  K\langle x\rangle
 \end{equation*}
 as the set of analytic functions on $\overline{B}(0,1)$. What about $\overline{B}(0,r)$ for $0<r<1$?
 By noting that $x\leq r\Leftrightarrow {\frac x r}\leq 1$, we have,
 \begin{equation*}
  K\left\langle {\frac x r}\right\rangle = \left\{\sum a_i x^i\sep \vert a_ir^i\vert\to 0\right\}.
 \end{equation*}
 is such set of analytic functions. The analytic functions on $B(0,1)$ is defined,
 \begin{equation*}
  \bigcap_{0<r<1}K\left\langle {\frac x r}\right\rangle = \left\{\sum a_ix^i \sep \vert a_ir^i\vert\to0,\forall 0<r<1\right\}.
 \end{equation*}
 \begin{proposition}[Universal Property of Tate algebra]
  For any complete $K$-algebra $A$ with $(a_1,\cdots,a_n)$ ordered set of n pb elements in $A$, there s a unique continuous morphism of $K$-algebras,
  \begin{equation*}
    f:\T_n\to A
  \end{equation*}
  such that $f(x_i)=a_i$.
 \end{proposition}
 \begin{proposition}
  Every filter is contained in a maximal filter.
 \end{proposition}
 \begin{proof}
  Consider a sequence of filters, 
  \begin{equation*}
    \mathcal{U}_0\subseteq\cdots
  \end{equation*}
  Then we claim,
  \begin{equation*}
    \mathcal{U} = \bigcup_{i\in I}\mathcal{U}_i
  \end{equation*}
  is again a filter. Indeed if $M,N\in\mathcal{U}$, we have,
  \begin{equation*}
    M\cap N\in\mathcal{U}_l\subseteq\mathcal{U}.
  \end{equation*}
  Also $\emptyset\not\in\mathcal{U}$.  If $M\in\mathcal{U}$ and $M\subseteq N$ then take $M\in\mathcal{U}$ then $N\in\mathcal{U}_l\subseteq \mathcal{U}$. By Zorn's lennma, every filter is contained in an ultrafilter.
 \end{proof}
 \begin{proposition}
  The direct image of an ultrafilter is an ultrafilter. If $\iota:A\hookrightarrow X$ then $\iota_*$ is a bijection from ultrafilters on $A$ to ultrafilters on $X$ containing $A$.
 \end{proposition}
 \begin{proof}
  Recall that the following statements are equivalent.
  \begin{enumerate}
    \item $\mathcal{U}$ is an ultrafilter on $X$
    \item For any subset $M\subseteq X$, either $M\in X$ or $X\backslash M\in X$. 
    \item If $M\cup N\in\mathcal{U}$ then either $M\in\mathcal{U}$ or $N\in\mathcal{U}$. 
  \end{enumerate}
  Suppose $f:X\to Y$ and $\mathcal{U}$ is an ultrafilter on $X$, then define,
  \begin{equation*}
    f_*\mathcal{U} = \{V\subseteq Y\sep f^{-1}(V)\in\mathcal{U}\}.
  \end{equation*}
  Easy to veryfy this is again an ultrafilter. For $M\subseteq Y$, note that 
  \begin{equation*}
    f^{-1}(M)\sqcup f^{-1}(Y\backslash M)
  \end{equation*}
  thus exactly one of $f^{-1}(M)$ and $f^{-1}(Y\backslash M)$ is in $\mathcal{U}$ hence one of $M$ and $Y\backslash M$ is in $f_*\mathcal{U}$. 
 \end{proof}
 \section{Week 2}
  \begin{lemma}
  Let $\mathcal{B}$ be an open basis of the space $X$. If $X$ is $T_0$ then $X$ is spectral if and only if every ultrafilter $\U$ has a generic limit $g$ such that $g\in \Omega\Leftrightarrow \Omega\in \U$.
  If this holds then for any $\mathcal{B}\ni\Omega$ is quasi-compact.
 \end{lemma}
 \par\textbf{Problem 5}\\
 We already have seen $\Spec A$ is spectral. It is $T_0$ and for any ultrafilter $\U$ on $\Spec A$, define,
 \begin{equation*}
  \g_\U = \{x\in A\sep \{\p \in\Spec A\sep x\in \p\}\in\U\}
 \end{equation*}
 is an ideal. This is quite obvious as $\p$ is an ideal. That is 
 \begin{equation*}
  \{\p\sep ax\in\p\}\supseteq\{\p\sep x\in \p\}.
 \end{equation*}
 \par\textbf{Problem 6}\\
 First prove that $X=\Sper A$ is $T_0$. If $P\not=Q$ then there is $a\in P\backslash Q$. Thus $P\not\in\mathcal{P}(-a)$ but $Q\in\mathcal{P}(-a)$. Let $\U$ be an ultrafilter on $X$. We then define,
 \begin{equation*}
  \mathfrak{P} = \{a\in A\sep \{C\in\Sper A\sep a\in C\}\in\U\}.
 \end{equation*}

 We would like to verify that $\mathfrak{P}$ is a generic limit. Let,
 \begin{equation*}
  F(a) \defeq \{\p\in\Sper A\sep a\in\p\}.
 \end{equation*} 
 For a square $b^2$, we have $F(b^2) = X$. Thus $b^2\in \mathfrak{P}$. For $F(-1) = \emptyset\not\in\U$ thus $-1\not\in \mathfrak{P}$.
 \par Let $a,b\in\mathfrak{P}$ then $F(a+b)\supseteq F(a)\cap F(b)$. Thus $F(a+b)\in\U$. Also for $F(ab)$. If $ab\in\mathfrak{P}$ but $a\not\in \mathfrak{P}$, then 
 \begin{equation*}
  \{\p\sep -b\in\p\}\supseteq\{\p\sep ab\in\p,a\not\in\p\} = F(ab)\cap (X\backslash F(a)).
 \end{equation*}
Since $F(a)$ is not in $\U$ thus $X\backslash F(a)$ is in $\U$. Thus $F(-b)\in\U$. 
\par For $\mathcal{P}(a)\in\mathcal{B}$, $\mathfrak{P}\in\mathcal{P}(a)$ if and only if $-a\not\in \mathfrak{P}$. Thus $F(-a)\not\in \U$. By definition $F(-a)\cup \mathcal{P}(a)=X$. Thus $\mathcal{P}(a)\in\U$.
Also $\mathfrak{P}\in\mathcal{P}(a_1,\cdots,a_n)$ then $\mathfrak{P}\in\mathcal{P}(a_i)$ for all $i$. Thus $\mathcal{P}(a_i)\in\U$. In particular $\mathcal{P}(a_1,\cdots,a_n)\in\U$. By the lemma $X$ is spectral and all elements in $\mathcal{B}$ are quasi-compact.
\par\textbf{Prroblem 7}\\
Clearly $X$ is $T_0$ since for any $x\not=y\in X$. Without loss of generality, we assume $x\not=\eta$. Then $X\backslash\{x\}$ does the work. Let $\U$ be an ultrafilter on $X$. If $U\in\U$, then $X\backslash U$ is finite thus there is $\{x_i\}\in \U$. If $U$ is cofinite, then $x_i\in U\Leftrightarrow U\in\U$. So $U$ is quasi-compact. If $\U$ contains all non-empty open sets then $\eta$ the generic limit. So $X$ is spectral and cofinite subsets are quasi-compact.
\par\textbf{Problem 8}\\
Note that if $x\not =\eta$, then $\{x\}$ is open thus this is a open neighborhood. On the other hand, the neighbourhood basis of $\eta$ is the basis open basis is the neighborhood basis as $X\backslash U_i\ni \eta$ then $U_i=\emptyset$ thus it is the whole.
\begin{theorem}
  If $X$ is quasi-compact and Hausdorff and denote $C(X)$ be the set of continuous functions $X\to \R$. Then $C(X)$ is a commutative unital ring. Then there is a one-to-one correspondence between $X$ and the maximal ideals of $C(X)$, mapping $x\mapsto \m_x = \{f|f(x)=0\}$.   
\end{theorem}
Are there non-maximal prime ideal? The answer is if $X$ is first countable, then non-maximal prime ideal exists if and only if $X$ is an infinite set.
\end{document}